\documentclass[11pt,a4paper]{article}

% --- Encoding & language ---
\usepackage[T1]{fontenc}
\usepackage[utf8]{inputenc}
\usepackage[italian]{babel}

% --- Typography ---
\usepackage{lmodern}
\usepackage{microtype}
\usepackage{geometry}
\geometry{margin=2.5cm, top=3cm, bottom=3cm}

% --- Math ---
\usepackage{amsmath,amssymb}

% --- Tables ---
\usepackage{booktabs}
\usepackage{longtable}
\usepackage{array}
\usepackage{tabularx}

% --- Graphics & colour ---
\usepackage{xcolor}
\definecolor{uniblu}{RGB}{0,119,200}
\definecolor{lightgray}{RGB}{245,245,245}
\definecolor{darkgray}{RGB}{60,60,60}

% --- Hyperlinks ---
\usepackage[colorlinks=true, linkcolor=uniblu, urlcolor=uniblu, citecolor=uniblu]{hyperref}

% --- Code listings ---
\usepackage{listings}
\lstset{
  basicstyle=\ttfamily\small,
  backgroundcolor=\color{lightgray},
  frame=single,
  framerule=0.4pt,
  rulecolor=\color{gray},
  breaklines=true,
  breakatwhitespace=false,
  columns=fullflexible,
  keepspaces=true,
  xleftmargin=1em,
  xrightmargin=1em,
  aboveskip=0.5em,
  belowskip=0.5em,
}

% --- Section styling ---
\usepackage{titlesec}
\titleformat{\section}[block]
  {\large\bfseries\color{uniblu}}{\thesection}{0.8em}{}[\titlerule]
\titleformat{\subsection}[block]
  {\normalsize\bfseries\color{darkgray}}{\thesubsection}{0.6em}{}
\titleformat{\subsubsection}[block]
  {\normalsize\itshape\color{darkgray}}{\thesubsubsection}{0.6em}{}

% --- Header / Footer ---
\usepackage{fancyhdr}
\pagestyle{fancy}
\fancyhf{}
\fancyhead[L]{\small\color{gray}BridgeManager 2025 --- Rapporto Tecnico del Modello}
\fancyhead[R]{\small\color{gray}\thepage}
\fancyfoot[C]{\small\color{gray}Giulio Mariniello --- Universit\`a degli Studi di Napoli Federico II}
\renewcommand{\headrulewidth}{0.4pt}
\renewcommand{\footrulewidth}{0.4pt}

% --- Misc ---
\usepackage{parskip}
\setlength{\parskip}{6pt}
\usepackage{enumitem}
\setlist{itemsep=2pt, topsep=4pt}

% ======================================================
\begin{document}

% --- Title page ---
\begin{titlepage}
  \centering
  \vspace*{2cm}
  {\color{uniblu}\rule{\linewidth}{2pt}}\\[1em]
  {\Huge\bfseries BridgeManager 2025}\\[0.4em]
  {\Large Rapporto Tecnico del Modello di Simulazione}\\[1em]
  {\color{uniblu}\rule{\linewidth}{1pt}}\\[2em]
  {\large Giulio Mariniello}\\[0.3em]
  {\normalsize Universit\`a degli Studi di Napoli Federico II\\
  Dipartimento di Strutture per l'Ingegneria e l'Architettura}\\[3em]
  {\normalsize
    \textbf{Versione:} 2025\\[0.3em]
    \textbf{Data:} 16 maggio 2026\\[0.3em]
    \textbf{Piattaforma:} Streamlit + Docker su Hugging Face Spaces\\[0.3em]
    \textbf{Riferimenti normativi principali:} EN 1990, EN 1991-2, JCSS PMC (2001), NTC 2018,\\
    MIT Linee Guida Ponti (2020)
  }\\[3em]
  \vfill
  {\small\color{gray}Documento generato automaticamente da Claude Code --- Anthropic}
\end{titlepage}

% --- ToC ---
\tableofcontents
\newpage

% ======================================================
\section{Architettura Generale}

Il simulatore modella la gestione di un portafoglio di ponti su un orizzonte temporale discretizzato in \textbf{trimestri} (un passo = 3 mesi). Ogni trimestre, un agente (umano o algoritmico) osserva lo stato dei propri ponti e pianifica interventi. L'ambiente avanza deterministicamente la fisica del degrado, applica le azioni pianificate, aggiorna l'affidabilit\`a strutturale, calcola i flussi economici e verifica le condizioni di fine-partita.

La sequenza di operazioni per ogni turno $t$ \`e:

\begin{enumerate}
  \item Completamento interventi scaduti (applicazione dei \textit{gain} ai componenti)
  \item Raccolta dei piani degli agenti
  \item Esecuzione delle azioni (analisi / monitoraggio / interventi / proof load)
  \item Degrado stocastico dei componenti non in manutenzione
  \item Aggiornamento previsioni di degrado per i ponti in monitoraggio attivo
  \item Evento sismico stocastico (se abilitato)
  \item Maturazione interessi sul budget + incassi trimestrali
  \item Verifica condizioni di fallimento (crollo strutturale o bancarotta)
  \item Aggiornamento affidabilit\`a di rete $\rightarrow$ \texttt{quarter} $+= 1$
\end{enumerate}

I dati di stato vengono serializzati in \texttt{game\_state.json} ad ogni turno; i log storici (azioni, budget, KPI) vengono aggiunti ai rispettivi file CSV di sessione.

% ======================================================
\section{Modello di Degrado dei Componenti}

\subsection{Struttura di un Componente}

Ogni ponte \`e composto da 4 \textbf{componenti strutturali}: \texttt{impalcato}, \texttt{pile}, \texttt{spalle}, \texttt{appoggi}. Ogni componente \`e descritto dai seguenti attributi:

\begin{center}
\begin{tabular}{lp{9cm}}
\toprule
\textbf{Attributo} & \textbf{Significato} \\
\midrule
\texttt{current\_condition} $\in [0,1]$ & Condizione strutturale reale (0~=~collasso, 1~=~perfetto) \\
\texttt{effective\_age\_years} & Et\`a efficace (aumenta in base a traffico, ambiente, stocasticit\`a) \\
\texttt{vita\_utile\_anni} & Vita utile nominale del componente (variabile casuale $\sim \mathcal{N}(L,\; 0.10\,L)$) \\
\texttt{epistemic\_uncertainty} $\in [0.03, 0.40]$ & Incertezza epistemica sulla resistenza \\
\texttt{observed\_condition\_min/max} & Intervallo di confidenza post-ispezione (\texttt{None} se mai ispezionato) \\
\bottomrule
\end{tabular}
\end{center}

\subsection{Legge di Degrado (Funzione di Sopravvivenza Weibull)}

La condizione di un componente all'et\`a efficace $t$ \`e:

\begin{equation}
  C(t) = \exp\!\left[-\lambda \left(\frac{t}{L}\right)^k\right]
\end{equation}

dove:
\begin{itemize}
  \item $L$ = vita utile del componente (anni)
  \item $k = \text{exponent} \times \text{shape\_factor}$ (esponente di Weibull, tipicamente $k \approx 3.0 \times 1.05 \div 1.35$)
  \item $\lambda$ = scala di degrado, campionata come $\lambda \sim \mathcal{N}(\mu_\lambda, 0.10)$, con $\lambda \geq 0.30$
\end{itemize}

I parametri nominali per ciascun componente (\texttt{COMPONENT\_PROFILES}) sono:

\begin{center}
\begin{tabular}{lcccc}
\toprule
\textbf{Componente} & \textbf{shape\_factor} & \textbf{scale\_mean} & \textbf{traffic\_sens} & \textbf{env\_sens} \\
\midrule
impalcato & 1.05 & 0.95 & 0.95 & 0.85 \\
pile      & 1.20 & 1.10 & 0.70 & 0.95 \\
spalle    & 1.10 & 1.00 & 0.60 & 1.00 \\
appoggi   & 1.35 & 1.12 & 0.75 & 1.10 \\
\bottomrule
\end{tabular}
\end{center}

Tutti i parametri sono campionati da distribuzioni normali con la media del profilo e $\sigma \approx 0.05\div0.10$, garantendo eterogeneità tra ponti diversi.

\subsection{Avanzamento dell'Età Efficace per Turno}

Ogni trimestre, per i componenti non in manutenzione attiva, l'età efficace avanza di:

\begin{equation}
  \Delta t_{\text{eff}} = 0.25 \cdot \tau_{\text{traf}}^{\alpha_{\text{traf}}}
                         \cdot E_{\text{env}}^{\alpha_{\text{env}}}
                         \cdot P_{\text{proc}}
\end{equation}

dove:
\begin{itemize}
  \item $\tau_{\text{traf}} = 1 + k_{\text{traf}} \cdot \dfrac{\text{ADT}}{\text{ADT}_{\text{ref}}}$ \quad (moltiplicatore traffico; default $k_{\text{traf}} = 0.05$, $\text{ADT}_{\text{ref}} = 10\,000$ veh/giorno)
  \item $E_{\text{env}} \sim \mathcal{N}(1.0,\; 0.10)$ \quad (fattore ambientale stocastico, campionato una volta per turno, condiviso da tutti i ponti)
  \item $P_{\text{proc}} \sim \text{LogN}(0,\; \sigma_{\text{proc}})$ \quad (processo stocastico per componente; $\sigma_{\text{proc}} = 0.12$, clampato $\geq 0.35$)
  \item $\alpha_{\text{traf}}$, $\alpha_{\text{env}}$ = \texttt{traffic\_sensitivity} e \texttt{environment\_sensitivity} del componente
\end{itemize}

Questo approccio cattura: (a) l'accelerazione del degrado per alto traffico, (b) la correlazione parziale tra ponti tramite il fattore ambientale condiviso, (c) la variabilità idiosincratica per componente tramite il moltiplicatore lognormale di processo.

\subsection{Aggiornamento della Condizione dopo Intervento}

Quando un intervento termina, ogni componente target riceve un \textit{gain} $g \sim \mathcal{N}(g_{\text{nom}},\; \sigma_g)$, clampato a $g \geq 0$:

\begin{itemize}
  \item Se $g \geq 1.0$ (sostituzione completa): $C \leftarrow 1.0$, $t_{\text{eff}} \leftarrow 0$
  \item Altrimenti: $C \leftarrow \min(1.0,\; C + g)$, poi inversione della legge di Weibull:
  \begin{equation}
    t_{\text{eff}} = L \cdot \left(\frac{-\ln C}{\lambda}\right)^{1/k}
  \end{equation}
\end{itemize}

Dopo ogni intervento: $\varepsilon \leftarrow \max(0.06,\; \varepsilon \times 0.90)$ (riduzione del 10\% dell'incertezza epistemica).

\subsection{Aumento dell'Incertezza Epistemica con il Tempo}

Ogni trimestre senza ispezione:
\begin{equation}
  \varepsilon \leftarrow \min(0.40,\; \varepsilon \times 1.01)
\end{equation}
Simula il progressivo deterioramento della conoscenza in assenza di monitoraggio.

% ======================================================
\section{Modello di Affidabilità Strutturale}

\subsection{Gerarchia di Valutazione}

Per ogni componente vengono valutati \textbf{3 stati limite}:

\begin{center}
\begin{tabular}{lcl}
\toprule
\textbf{Stato Limite} & \textbf{Categoria} & \textbf{Modalità di Rottura} \\
\midrule
Flessione (\textit{flexure}) & ULS & Duttile \\
Taglio (\textit{shear})      & ULS & Fragile \\
Fessurazione (\textit{cracking}) & SLS & Servizio \\
\bottomrule
\end{tabular}
\end{center}

I risultati ULS di ogni componente vengono combinati in \textbf{serie} per ottenere l'affidabilità del componente. I risultati di tutti i componenti strutturali del ponte vengono a loro volta combinati in serie per ottenere l'affidabilità del ponte.

\subsection{Variabili Aleatorie}

Per ogni stato limite vengono campionati $N$ (default $N=256$) vettori di variabili aleatorie:

\begin{align}
  R_i &= R_{\mu} \cdot \xi_{\text{mat}} \cdot \xi_{\text{geom}} \cdot \theta_R \\
  S_i &= \theta_S \cdot (G_i + Q_i)
\end{align}

dove:
\begin{itemize}
  \item $\xi_{\text{mat}} \sim \text{LogN}(1, v_m)$ --- incertezza sul materiale
  \item $\xi_{\text{geom}} \sim \text{LogN}(1, v_g)$ --- incertezza geometrica
  \item $\theta_R \sim \text{LogN}(1, v_{\text{mod},R})$ --- incertezza di modello sulla resistenza
  \item $\theta_S \sim \text{LogN}(1, v_{\text{mod},S})$ --- incertezza di modello sulla domanda
  \item $G \sim \mathcal{N}(\mu_G, v_G)$ --- carico permanente (normale)
  \item $Q \sim \text{Gumbel}(\mu_Q, v_Q)$ --- carico variabile (Gumbel)
\end{itemize}

\subsection{Resistenza Media}

\begin{equation}
  R_{\mu} = \xi_{\text{bias}} \cdot \max(C,\; 0.02) \cdot \bigl(2.24 + 0.20 \cdot (1 - \varepsilon)\bigr)
\end{equation}

Il fattore $2.24$ \`e il bias nominale (da calibrare su dati AINOP/JCSS). Le varianti per stato limite sono:
\begin{itemize}
  \item Flessione: $R_{\mu} \times 1.08$
  \item Taglio: $R_{\mu} \times (0.92 + 0.06 \cdot \alpha_{\text{traf}})$
  \item Fessurazione: $R_{\mu} \times \xi_{\text{crack}} \times 0.76$
\end{itemize}

\subsection{COV delle Distribuzioni}

\begin{center}
\begin{tabular}{ll}
\toprule
\textbf{Parametro} & \textbf{Formula} \\
\midrule
$v_m$ (materiale) & $\min(0.24,\; 0.04 + 0.40 v_{\text{res}} + 0.05(1-C))$ \\
$v_g$ (geometria) & $\min(0.10,\; v_{g,\text{base}} \cdot (1 + 0.12 \cdot k_{\text{shape}}))$ \\
$v_{\text{mod},R}$ (modello resistenza) & $\min(0.16,\; v_{\text{mod}} \cdot (1 + 0.75\varepsilon))$ \\
$v_{\text{mod},S}$ (modello domanda) & $\min(0.14,\; 0.04 + 0.30 v_{\text{demand}} + 0.05\varepsilon)$ \\
$v_G$ (carico permanente) & $\min(0.08,\; \max(0.05,\; v_{\text{demand}} \cdot 0.50))$ \\
$v_Q$ (carico variabile) & $\min(0.18,\; \max(v_{\text{demand}},\; 0.04 + 0.22\beta_{\eta} + 0.06(1-C)))$ \\
\bottomrule
\end{tabular}
\end{center}

\subsection{Stima FOSM dai Momenti Campionari}

Per evitare l'instabilità del conteggio diretto dei fallimenti a $N$ piccolo (che restituisce 0 eventi per $\beta > 3$ nel 97\% dei campioni con $N=256$), si usa la formula \textbf{FOSM lognormale} applicata ai momenti campionari di $R$ e $S$:

\begin{align}
  \sigma_{\ln R} &= \sqrt{\ln(1 + \text{COV}_R^2)} \\
  \sigma_{\ln S} &= \sqrt{\ln(1 + \text{COV}_S^2)} \\
  \beta_{\text{ann}} &= \frac{\ln(\mu_R / \mu_S)}{\sqrt{\sigma_{\ln R}^2 + \sigma_{\ln S}^2}}
\end{align}

Questa formulazione \`e equivalente all'approccio FOSM standard (Hasofer-Lind per variabili lognormali) ma usa momenti stimati dai campioni, garantendo stabilit\`a per qualsiasi valore di $\beta$ e $N$.

\subsection{Probabilità di Fallimento e Periodo di Riferimento}

Per un singolo stato limite con $\beta_{\text{ann}}$ annuale:
\begin{equation}
  p_{f,\text{ann}} = \Phi(-\beta_{\text{ann}})
\end{equation}

La probabilità cumulata su periodo di riferimento $T$ (stazionaria):
\begin{equation}
  p_{f,T} = 1 - (1 - p_{f,\text{ann}})^T
  \qquad
  \beta_T = -\Phi^{-1}(p_{f,T})
\end{equation}

Per la valutazione \textbf{non-stazionaria} (affidabilità del ciclo di vita), si proietta il profilo anno per anno per i prossimi $T_{\text{proj}}$ anni e si calcola:
\begin{equation}
  p_{f,TV} = 1 - \prod_{t=1}^{T_{\text{proj}}} \bigl[1 - p_{f,\text{ann}}(t)\bigr]
\end{equation}
in accordo con Rackwitz (2001) e Frangopol et al.\ (2004). Questo valore viene usato come \texttt{bridge.failure\_probability}.

\subsection{Target di Affidabilità (JCSS PMC / EN 1990)}

\begin{center}
\begin{tabular}{lccc}
\toprule
\textbf{Classe di Conseguenza} & $\beta_{\text{rif}}$ ($T=50$ anni) & $\beta_{\text{ann}}$ (derivato) & $p_{f,\text{ann}}$ \\
\midrule
CC1 & 3.3 & 4.2 & $1.3 \times 10^{-5}$ \\
CC2 (default) & 3.8 & 4.7 & $1.3 \times 10^{-6}$ \\
CC3 & 4.3 & 5.2 & $1.0 \times 10^{-7}$ \\
\midrule
SLS (tutti) & 1.5 & --- & --- \\
ULS fragile ($+\Delta\beta$) & $\beta_{\text{CC}} + 0.30$ & --- & --- \\
\bottomrule
\end{tabular}
\end{center}

\subsection{Combinazione a Sistema in Serie}

Sotto l'ipotesi di indipendenza (conservativa per componenti positivamente correlati):

\begin{equation}
  p_{f,\text{sistema}} = 1 - \prod_i (1 - p_{f,i})
\end{equation}

Applicata prima ai 3 stati limite del componente, poi ai componenti strutturali del ponte, con l'eventuale aggiunta della componente sismica \texttt{extra\_annual\_pf}.

% ======================================================
\section{Modello Informativo: Ispezioni e Incertezza Epistemica}

\subsection{Osservazione della Condizione}

Un'ispezione genera un'osservazione rumorosa della condizione reale:

\begin{equation}
  \sigma_{\text{obs}} = \max(0.005,\; \sigma_{\text{base}} \cdot (0.65 + \varepsilon))
\end{equation}

L'intervallo di confidenza osservato \`e:
\begin{equation}
  [C_{\text{obs,min}},\; C_{\text{obs,max}}] = \bigl[\hat{C} - 1.96\,\sigma_{\text{obs}},\; \hat{C} + 1.96\,\sigma_{\text{obs}}\bigr] \cap [0,1]
\end{equation}

Il valore di $\sigma_{\text{base}}$ per ciascuna azione \`e:

\begin{center}
\begin{tabular}{lc}
\toprule
\textbf{Azione} & $\sigma_{\text{base}}$ \\
\midrule
L1\_inspect & 0.10 \\
L4\_verify  & 0.01 \\
monitoring  & 0.005 \\
\bottomrule
\end{tabular}
\end{center}

\subsection{Aggiornamento dell'Incertezza Epistemica}

Ogni ispezione riduce $\varepsilon$ tramite il \texttt{knowledge\_gain} $\kappa$ dell'azione:
\begin{equation}
  \varepsilon_{\text{new}} = \max(0.03,\; \varepsilon \cdot (1 - \kappa))
\end{equation}

Valori: monitoring $\kappa = 0.30$; proof\_load\_test $\kappa = 0.50$.

\subsection{Classe di Attenzione (CdA)}

Calcolata sulla condizione osservata minima tra tutti i componenti ispezionati:

\begin{center}
\begin{tabular}{lcl}
\toprule
\textbf{CdA} & \textbf{Soglia osservata} & \textbf{Scadenza ispezione} \\
\midrule
Bassa       & $> 0.90$ & $+8$ trimestri \\
Medio-Bassa & $> 0.80$ & $+6$ trimestri \\
Media       & $> 0.60$ & $+4$ trimestri \\
Medio-Alta  & $> 0.40$ & $+2$ trimestri \\
Alta        & $\leq 0.40$ & $+1$ trimestre \\
Sconosciuta & mai ispezionato & --- \\
\bottomrule
\end{tabular}
\end{center}

Con \texttt{enforce\_guidelines = True}, il mancato rispetto delle scadenze genera un \textit{backlog} conteggiato nei KPI.

\subsection{Monitoraggio Continuo}

Quando il monitoraggio \`e attivo, il simulatore calcola ad ogni turno la curva predittiva di degrado per i prossimi $\texttt{monitoring\_duration\_years} \times 4$ trimestri, proiettando l'et\`a efficace con il moltiplicatore medio di traffico e ambiente (componente deterministica). La banda di confidenza \`e $\pm 1.96\,\sigma_{\text{obs}}$ con $\sigma_{\text{base}} = 0.005$.

% ======================================================
\section{Prova di Carico (Proof Load Test)}

\subsection{Logica del Test a Passi}

Il test viene eseguito a \texttt{step\_count} (default 3) livelli di carico crescenti $\alpha_k$, da $\alpha/\text{step}$ ad $\alpha$:

\begin{equation}
  S_{\text{test},k} = \theta_{\text{test}} \cdot (G + \alpha_k \cdot Q_{\text{codice}})
\end{equation}

Per ogni campione $i$ e passo $k$:
\begin{itemize}
  \item $R_i \leq S_{\text{test},k}$ su stato SLS $\Rightarrow$ campione \textbf{schermato} (fermato preventivamente)
  \item $R_i \leq S_{\text{test},k}$ su stato ULS $\Rightarrow$ campione \textbf{fallito} (collasso durante test)
\end{itemize}

Al termine del test:
\begin{align}
  p_{\text{pass}} &= P(\text{sopravvissuto a tutti i passi}) \\
  p_{\text{screen}} &= P(\text{schermato senza collasso}) \\
  p_{\text{fail,test}} &= P(\text{collasso durante il test})
\end{align}

\subsection{Aggiornamento Bayesiano Analitico}

Per ogni stato limite, si condiziona la resistenza $R \sim \text{LogN}$ sull'evento $R > x_0$, con soglia:
\begin{equation}
  x_0 = \mu_G + \alpha \cdot Q_{\text{codice}}
\end{equation}

I momenti esatti della distribuzione troncata (Aitchison \& Brown 1957; JCSS PMC \S3.7) sono:

\begin{align}
  k &= \frac{\ln x_0 - \mu_{\ln R}}{\sigma_{\ln R}} \\
  \Phi_0 &= \Phi(-k) = P(R > x_0) \\
  \mathbb{E}[R \mid R > x_0] &= R_\mu \cdot \frac{\Phi(\sigma_{\ln R} - k)}{\Phi_0} \\
  \mathbb{E}[R^2 \mid R > x_0] &= R_\mu^2 \cdot e^{\sigma_{\ln R}^2} \cdot \frac{\Phi(2\sigma_{\ln R} - k)}{\Phi_0} \\
  \text{COV}_{\text{post}} &= \frac{\sqrt{\mathbb{E}[R^2] - \mathbb{E}[R]^2}}{\mathbb{E}[R]}
\end{align}

La decomposizione dell'incertezza preserva le componenti irreducibili:
\begin{equation}
  v_{m,\text{post}} = \sqrt{\max\!\bigl(0,\; \text{COV}_{\text{post}}^2 - v_g^2 - v_{\text{mod}}^2\bigr)}
\end{equation}

\textbf{Propriet\`a formale:} il condizionamento su $R > x_0$ implica $\mathbb{E}[R\mid R>x_0] \geq R_\mu$, quindi $\beta_{\text{post}} \geq \beta_{\text{prior}}$ per costruzione. Il proof load test non pu\`o peggiorare la stima di affidabilit\`a di un ponte che lo supera.

% ======================================================
\section{Modello di Fragilità al Traffico}

\subsection{Modelli di Carico Disponibili}

\begin{center}
\begin{tabular}{lccc}
\toprule
\textbf{Modello} & $\eta$ (mediana) & $\beta$ (dispersione) & \textbf{Uso tipico} \\
\midrule
GL\_light  & 7.17 & 0.40 & Traffico leggero \\
GL\_medium & 3.19 & 0.36 & Traffico misto \\
GL\_heavy  & 2.33 & 0.33 & Traffico pesante (default) \\
NTC        & 1.68 & 0.33 & NTC 2018 LM1 \\
\bottomrule
\end{tabular}
\end{center}

\subsection{Fattori di Correzione}

\begin{equation}
  \eta_{\text{eff}} = \eta_{\text{base}}
    \cdot \gamma_{\text{corsie}}
    \cdot \gamma_{\text{offset}}
    \cdot \gamma_{\text{area}}
    \cdot \gamma_{\text{age}}
\end{equation}

\begin{itemize}
  \item $\gamma_{\text{corsie}}$: $+10\%$ per 2 corsie, $+35\%$ per 1 corsia
  \item $\gamma_{\text{offset}} = 1 + 0.22 \cdot d_{\text{off}}$ --- beneficio dell'eccentricit\`a del carico
  \item $\gamma_{\text{area}} = 1.03$ se l'area rimanente \`e inattiva
  \item $\gamma_{\text{age}} = 1 - 0.18 \cdot (1 - C_{\text{str,min}})$ --- penale per degrado strutturale
\end{itemize}

dove $C_{\text{str,min}} = \min(C_{\text{imp}}, C_{\text{pile}}, C_{\text{spalle}})$ \`e la condizione minima dei soli componenti portanti.

Il tasso di fallimento annuale al traffico \`e:
\begin{equation}
  \lambda_{\text{ann}} = \begin{cases}
    p_{f,\alpha=1} \times P_{\text{overload,month}} \times P_{\text{congestion}} & \text{(GL\_heavy)} \\
    p_{f,\alpha=1} \times 10^{-3} & \text{(NTC)}
  \end{cases}
\end{equation}

% ======================================================
\section{Modello Sismico}

Il modulo sismico usa un \textbf{modello di fragilità lognormale} calibrato su classi di ponti e stati limite (GdL, SD, NC):

\begin{equation}
  P(\text{danno} \mid \text{PGA}) = \Phi\!\left(\frac{\ln(\text{PGA}/\eta_s)}{\beta_s}\right)
\end{equation}

Il PGA controllante viene stimato empiricamente come:
\begin{equation}
  \text{PGA}_{\text{ctrl}} = 0.04 + 0.18 \cdot p_{f,\text{ann}} + 0.015 \cdot M_{\text{max}}
\end{equation}

La probabilit\`a annuale di eccedenza sismica \`e:
\begin{equation}
  P_{\text{sism,ann}} = P(\text{danno} \mid \text{PGA}_{\text{ctrl}}) \times P(\text{terremoto})
\end{equation}

Gli eventi sismici stocastici vengono generati con:
\begin{itemize}
  \item Probabilit\`a trimestrale: $p_{\text{EQ}} = p_{\text{ann}} / 4$
  \item Magnitudo $M \sim \mathcal{U}[M_{\min}, M_{\max}]$
  \item Posizione: casuale nel raggio $1.5\,r$ attorno alle coordinate del sito
\end{itemize}

% ======================================================
\section{Modello di Rete: Importanza e Priorità}

\subsection{Fattore di Importanza della Rete (NIF)}

\begin{equation}
  \text{NIF} = w_{\text{ADT}} \cdot \frac{\text{ADT}_i}{\text{ADT}_{\max}}
              + (1 - w_{\text{ADT}}) \cdot \frac{c_i}{c_{\max}}
\end{equation}

dove $c_i$ \`e la centralit\`a inversamente proporzionale alla distanza media haversine ai 3 ponti pi\`u vicini. Peso default: $w_{\text{ADT}} = 0.65$.

\subsection{Risk Priority Index (RPI)}

\begin{equation}
  \text{RPI} = \text{NIF} \times p_{f,\text{ann}}
\end{equation}

Consente di priorizzare interventi su ponti ad alta importanza strategica e alta probabilit\`a di fallimento.

% ======================================================
\section{Modello Economico}

\subsection{Flusso di Budget per Turno}

Il budget matura \textbf{interessi asimmetrici} prima dell'incasso:

\begin{equation}
  B_{t+1} = \begin{cases}
    B_t \cdot (1 + r_d)^{0.25} + I_t & \text{se } B_t \geq 0 \quad (r_d = 2\%\text{/anno}) \\
    B_t \cdot (1 + r_\delta)^{0.25} + I_t & \text{se } B_t < 0 \quad (r_\delta = 5\%\text{/anno})
  \end{cases}
\end{equation}

dove $I_t$ \`e l'incasso trimestrale netto (ricavi meno costi degli interventi avviati).

\subsection{Incasso Trimestrale}

Il moltiplicatore di ricavo dipende dall'Health Index $\text{HI}$ e dallo stato di manutenzione:

\begin{equation}
  \mu_{\text{rev}} = \begin{cases}
    1 - r_{\text{impact}}           & \text{se in manutenzione attiva} \\
    1.0                              & \text{se } \text{HI} > 0.5 \\
    \dfrac{\text{HI} - 0.2}{0.30}   & \text{se } 0.2 < \text{HI} \leq 0.5 \\
    0.0                              & \text{se } \text{HI} \leq 0.2
  \end{cases}
\end{equation}

\subsection{Health Index (HI)}

Media ponderata della condizione corrente dei componenti:

\begin{equation}
  \text{HI} = \frac{\sum_i w_i \cdot C_i}{\sum_i w_i}
\end{equation}

Pesi default: impalcato 50\%, pile 30\%, spalle 15\%, appoggi 5\%.

\subsection{Net Asset Value (NAV)}

\begin{equation}
  \text{NAV} = \underbrace{V_0 \cdot \text{HI}}_{\text{valore operativo}}
              - \underbrace{C_{\text{ripristino}} \cdot (1 - \text{HI})}_{\text{backlog manutentivo}}
              - \underbrace{I_{\text{trim}} \cdot r_{\text{impact}} \cdot q_{\text{res}}}_{\text{perdita ricavi attuale}}
\end{equation}

dove $V_0$ \`e il valore iniziale dell'asset, $C_{\text{ripristino}}$ il costo stimato di ripristino completo, $I_{\text{trim}}$ i ricavi trimestrali nominali, $r_{\text{impact}}$ l'impatto dell'intervento in corso, e $q_{\text{res}}$ i trimestri rimanenti dell'intervento. Il NAV \`e clampato a $\geq 0$.

% ======================================================
\section{Azioni di Manutenzione}

\begin{longtable}{p{3.5cm}p{1.5cm}p{2.5cm}p{1.5cm}p{4cm}}
\toprule
\textbf{Azione} & \textbf{Tipo} & \textbf{Costo} & \textbf{Durata} & \textbf{Effetto} \\
\midrule
\endhead
L1\_inspect & analysis & 1 €/m² & istantaneo & Osservazione $\pm10\%$ \\
L4\_verify & analysis & 100 €/m² & istantaneo & Osservazione $\pm1\%$ \\
monitoring & monitoring & 70 €/m² & istantaneo (attivo 5 anni) & Osservazione $\pm0.5\%$, previsioni future \\
proof\_load\_test & proof\_load & 25 €/m² & istantaneo & Aggiornamento Bayesiano di $\beta$ \\
replace\_impalcato & intervention & 4000 €/m² & 4 trimestri & gain = 1.0 su impalcato, 0.15 sugli altri \\
replace\_pile & intervention & 3000 €/m² & 3 trimestri & gain = 1.0 su pile, 0.15 sugli altri \\
replace\_spalle & intervention & 2500 €/m² & 3 trimestri & gain = 1.0 su spalle, 0.15 sugli altri \\
replace\_appoggi & intervention & €100\,000 fissi & 2 trimestri & gain = 1.0 su appoggi, 0.15 sugli altri \\
ripristino\_strutturale & intervention & 1500 €/m² & 3 trimestri & gain = 0.5 su pile + spalle ($\sigma_g = 0.05$) \\
ripristino\_completo & intervention & 3500 €/m² & 5 trimestri & gain: imp 0.7, pile 0.6, spalle 0.5, appoggi 0.4 ($\sigma_g = 0.05$) \\
\bottomrule
\end{longtable}

Il gain effettivo \`e $g_{\text{eff}} = \max(0,\; g_{\text{nom}} + \mathcal{N}(0, \sigma_g))$. Le sostituzioni complete (replace\_*) ripristinano il componente target a condizione 1.0 e forniscono un \textbf{bonus del 15\%} agli altri componenti per effetto trascinamento.

\subsection{Vincoli di Pianificazione}

\begin{itemize}
  \item Un ponte in manutenzione attiva non pu\`o ricevere nuove azioni.
  \item Non \`e possibile pianificare la stessa azione consecutivamente (eccezione: L1\_inspect).
  \item Non \`e possibile avviare il monitoraggio se gi\`a attivo.
  \item Il budget \`e verificato all'esecuzione: se insufficiente, l'azione \`e annullata.
  \item Le \texttt{allowed\_actions} per scenario filtrano le opzioni sin dall'avvio della sessione.
\end{itemize}

% ======================================================
\section{Condizioni di Fine-Partita}

\subsection{Crollo Strutturale}

L'agente viene eliminato se la \textit{condizione minima strutturale} di un qualsiasi ponte sotto la sua gestione scende sotto la soglia:

\begin{equation}
  C_{\text{str,min}} = \min(C_{\text{imp}}, C_{\text{pile}}, C_{\text{spalle}}) \leq 0.10
\end{equation}

Gli appoggi sono esclusi in quanto componenti funzionali, non strutturali portanti.

\subsection{Bancarotta}

L'agente viene eliminato se il budget rimane negativo per \textbf{4 trimestri consecutivi} (= 1 anno):
\begin{equation}
  \text{negative\_budget\_streak} \geq 4
\end{equation}

\subsection{Fine Naturale}

La simulazione termina al raggiungimento di \texttt{max\_quarters} o di \texttt{years} $\times 4$ trimestri.

% ======================================================
\section{Flusso Sperimentale}

La modalit\`a ricerca gestisce il ciclo di vita di ogni sessione attraverso 4 fasi:

\begin{enumerate}
  \item \textbf{tutorial} --- 5 passi didattici interattivi (ruolo, ponti, azioni, economia, consigli)
  \item \textbf{playing} --- gioco vero e proprio con dashboard interattiva (\texttt{max\_quarters} trimestri)
  \item \textbf{questionnaire} (opzionale) --- raccolta dati su esperienza, strategia, consapevolezza dell'asimmetria informativa
  \item \textbf{done} --- debrief con confronto contro baseline euristica eseguita post-hoc sullo stesso scenario
\end{enumerate}

Il questionario raccoglie: anni di esperienza, ruolo professionale, familiarit\`a con l'affidabilit\`a strutturale, strategia dichiarata, approccio sistematico vs.\ reattivo, consapevolezza dell'asimmetria informativa, uso dell'indice $\beta$ nelle decisioni.

% ======================================================
\section{KPI Registrati per Sessione}

\begin{center}
\begin{tabular}{lp{9cm}}
\toprule
\textbf{KPI} & \textbf{Calcolo} \\
\midrule
\texttt{health\_index} & $\sum(C_{\min,i} \cdot \text{ADT}_i) / \sum \text{ADT}_i$ (media pesata per traffico) \\
\texttt{total\_asset\_value} & $\sum V_{0,i} \cdot C_{\min,i}$ \\
\texttt{net\_worth} & budget $+$ total\_asset\_value \\
\texttt{backlog\_count} & ponti con scadenza ispezione superata \\
\texttt{inspection\_coverage} & ponti con almeno un'osservazione / totale \\
\texttt{knowledge\_gap\_count} & ponti mai ispezionati \\
\texttt{active\_monitoring\_count} & ponti con monitoraggio attivo \\
\texttt{network\_reliability\_index} & $\min(\beta_{\text{post}})$ su tutti i ponti \\
\texttt{network\_failure\_probability} & $\max(p_{f,\text{post}})$ su tutti i ponti \\
\bottomrule
\end{tabular}
\end{center}

Tutti i KPI vengono registrati in CSV ad ogni turno per l'analisi scientifica successiva.

% ======================================================
\section{Limitazioni e Assunzioni Note}

\begin{center}
\begin{tabular}{p{4cm}p{5.5cm}p{3cm}}
\toprule
\textbf{Limitazione} & \textbf{Impatto} & \textbf{Stato} \\
\midrule
Indipendenza componenti nel sistema serie & Conservativo: $p_{f,\text{sistema}}$ sovrastimata & Accettato per semplicità \\
Fattore 2.24 non calibrato su dati di campo & $R_\mu$ relativo, coerente internamente & Da calibrare su AINOP/JCSS \\
$\beta_{\text{post}} \geq \beta_{\text{prior}}$ per proof load & Propriet\`a matematica del condizionamento & Corretto e dichiarabile \\
Correlazione tra ponti solo via fattore ambientale & Correlazione strutturale non modellata & Accettato \\
Fragilità sismica empirica & Non sostituisce analisi sismica NTC & Scopo didattico \\
Revenue disconnesso da $\beta$ & --- & Design sperimentale intenzionale \\
\bottomrule
\end{tabular}
\end{center}

% --- Bibliography ---
\newpage
\begin{thebibliography}{99}
\bibitem{aitchison1957} Aitchison, J., Brown, J.A.C. (1957). \textit{The Lognormal Distribution}. Cambridge University Press.
\bibitem{jcss2001} JCSS (2001). \textit{Probabilistic Model Code}. Joint Committee on Structural Safety.
\bibitem{en1990} EN 1990 (2002). \textit{Eurocode --- Basis of Structural Design}. CEN.
\bibitem{en19912} EN 1991-2 (2003). \textit{Eurocode 1: Actions on structures --- Part 2: Traffic loads on bridges}. CEN.
\bibitem{ntc2018} NTC (2018). \textit{Norme Tecniche per le Costruzioni}. D.M. 17 gennaio 2018.
\bibitem{mit2020} MIT (2020). \textit{Linee Guida per la Classificazione e Gestione del Rischio, la Valutazione della Sicurezza e il Monitoraggio dei Ponti}. Ministero delle Infrastrutture e dei Trasporti.
\bibitem{rackwitz2001} Rackwitz, R. (2001). Reliability analysis --- a review and some perspectives. \textit{Structural Safety}, 23(4), 365--395.
\bibitem{frangopol2004} Frangopol, D.M., Kallen, M.-J., van Noortwijk, J.M. (2004). Probabilistic models for life-cycle performance of deteriorating structures. \textit{Progress in Structural Engineering and Materials}, 6(4), 197--212.
\bibitem{noortwijk2009} van Noortwijk, J.M. (2009). A survey of the application of gamma processes in maintenance. \textit{Reliability Engineering \& System Safety}, 94(1), 2--21.
\end{thebibliography}

\end{document}
